\documentclass[11pt]{article}
\usepackage[margin=1in]{geometry}
\usepackage{amsmath,amssymb,amsthm,mathtools,bm,booktabs,array,longtable,microtype}
\usepackage[hidelinks]{hyperref}
\usepackage{enumitem}
\newtheorem{theorem}{Theorem}[section]
\newtheorem{lemma}[theorem]{Lemma}
\newtheorem{proposition}[theorem]{Proposition}
\newtheorem{corollary}[theorem]{Corollary}
\theoremstyle{definition}\newtheorem{definition}[theorem]{Definition}
\theoremstyle{remark}\newtheorem{remark}[theorem]{Remark}
\title{A Finite Addressed Realization of the Standard-Model Gauge Representation on $V_7\otimes\mathbb C[\mathbb Z_3\times\mathrm{Dic}_6]$}
\author{UCD Mathematical Physics Program}
\date{September 2026}
\begin{document}
\maketitle

\begin{abstract}
Let
\[
\mathcal H_{504}=V_7\otimes\mathbb C[\mathbb Z_3]\otimes\mathbb C[\mathrm{Dic}_6],
\qquad \dim \mathcal H_{504}=504,
\]
where $V_7$ consists of one center state and a six-site ring, and
\[
\mathrm{Dic}_6=\langle a,x\mid a^{12}=1,\;x^2=a^6,\;xax^{-1}=a^{-1}\rangle.
\]
We construct commuting color, weak-isospin and hypercharge actions on this finite addressed space.  The six-leaf sector factors by the Chinese remainder theorem as $\mathbb C[\mathbb Z_6]\simeq\mathbb C^3\otimes\mathbb C^2$, yielding an $\mathfrak{su}(3)$ color action that annihilates the center.  The central subgroup $\langle a^6\rangle\simeq\mathbb Z_2$ yields a two-state weak fiber, while the one-dimensional character $\chi(a)=1$, $\chi(x)=-1$ supplies a chirality involution.  With $A$ denoting center/leaf parity, $B=L(a^6)$, and $P_L$ the chirality projector, the electric charge, weak generator and hypercharge are
\[
Q=\frac{B}{2}+\frac{A}{3}-\frac16 I,
\qquad
T_3=\frac12 P_LB,
\qquad
Y=2(Q-T_3)=\frac{2A-I}{3}+P_RB.
\]
After suppressing address multiplicities, the resulting isotypic sectors are exactly
\[
Q_L:(\mathbf3,\mathbf2)_{1/3},\quad
L_L:(\mathbf1,\mathbf2)_{-1},\quad
\nu_R:(\mathbf1,\mathbf1)_0,\quad
e_R:(\mathbf1,\mathbf1)_{-2},
\]
with the quark singlets $u_R:(\mathbf3,\mathbf1)_{4/3}$ and $d_R:(\mathbf3,\mathbf1)_{-2/3}$, and with the independent $\mathbb Z_3$ factor providing three family copies.  We verify the complete gauge Lie algebra, the Standard-Model electric charges and hypercharges including a right-handed neutrino, all one-generation gauge and mixed anomaly cancellations, the boundary trace ratio $G_Y/G_2=5/3$, and the one-loop coefficients $b_3=-7$, $b_2=-19/6$, $b_Y=41/6$.
\end{abstract}

\section{Finite addressed space}
Let $\{|c\rangle,|1\rangle,\ldots,|6\rangle\}$ be an orthonormal basis of $V_7$. Define
\[
P_c=|c\rangle\langle c|,\qquad P_\ell=I_7-P_c,
\qquad A=P_\ell-P_c=I_7-2P_c.
\]
Thus $A^2=I_7$ and the eigenvalue $A=-1$ labels the center while $A=+1$ labels the leaf sector.

On $\mathbb C[\mathrm{Dic}_6]$ let $L(g)$ be the left regular representation and set
\[
B=L(a^6).
\]
Because $a^6$ is central of order two,
\[
B^2=I_{24}.
\]
The lifted operators $A$ and $B$ commute on $\mathcal H_{504}$.

\begin{lemma}[Charge polynomial reduction]
Let $C=2A+B$. Then
\[
\frac{C^3-4C-3I}{18}
=\frac{B}{2}+\frac{A}{3}-\frac16I.
\]
\end{lemma}
\begin{proof}
Since $A^2=B^2=I$ and $[A,B]=0$,
\[
C^3=(2A+B)^3=8A+12B+6A+B=14A+13B.
\]
Hence $C^3-4C-3I=6A+9B-3I$, which gives the result.
\end{proof}

We therefore define
\begin{equation}
\label{eq:Q}
Q:=\frac{B}{2}+\frac{A}{3}-\frac16I.
\end{equation}
Its four $(A,B)$ eigenvalue pairs give
\[
(-1,-1)\mapsto -1,
\quad(-1,+1)\mapsto0,
\quad(+1,-1)\mapsto-\frac13,
\quad(+1,+1)\mapsto\frac23.
\]

\section{Color from the six-leaf ring}
The leaf space is the regular space of $\mathbb Z_6$.  The Chinese remainder map
\[
\mathbb Z_6\longrightarrow\mathbb Z_3\times\mathbb Z_2,
\qquad n\longmapsto(n\bmod3,n\bmod2)
\]
is a group isomorphism. Consequently
\[
P_\ell V_7\cong\mathbb C^3_{\rm col}\otimes\mathbb C^2_{\rm ori}.
\]
If $S_6$ is the unit shift on the ring, then in the CRT basis
\[
S_6=X_3\otimes X_2.
\]

Let $t_a=\lambda_a/2$, $a=1,\ldots,8$, be the Gell--Mann generators on $\mathbb C^3$, normalized by
\[
\operatorname{tr}(t_at_b)=\frac12\delta_{ab}.
\]
Define their action on $V_7$ by
\[
T_a^{(3)}=0\oplus(t_a\otimes I_2).
\]

\begin{theorem}[Native color algebra]
The operators $T_a^{(3)}$ satisfy the $\mathfrak{su}(3)$ commutation relations, annihilate the center, and obey
\[
\sum_{a=1}^8\left(T_a^{(3)}\right)^2=\frac43P_\ell.
\]
Thus the center is a color singlet and the leaf sector is a color triplet with a two-dimensional address multiplicity.
\end{theorem}
\begin{proof}
The result follows from the fundamental $\mathfrak{su}(3)$ representation on the first CRT factor and the identity action on the second.  The quadratic Casimir of the fundamental representation is $C_F=4/3$.
\end{proof}

\section{Weak fiber and chirality}
Let $Z=\langle a^6\rangle=\{1,a^6\}\simeq\mathbb Z_2$.  Choose a transversal $T$ for $\mathrm{Dic}_6/Z$.  Every group element is uniquely $tz^b$, $t\in T$, $b\in\{0,1\}$, and hence
\[
\mathbb C[\mathrm{Dic}_6]\cong\mathbb C^{12}\otimes\mathbb C^2_Z.
\]
In this factorization $B=L(a^6)=I_{12}\otimes\sigma_x$.  Let $Z_B=I_{12}\otimes\sigma_z$ and define
\[
\tau_1=Z_B,\qquad
\tau_2=-iBZ_B,\qquad
\tau_3=B.
\]
Then
\[
[\tau_i,\tau_j]=2i\epsilon_{ijk}\tau_k,
\qquad \tau_i^2=I.
\]
Different choices of transversal are related by a unitary change of the $Z$-fiber frame, so the resulting $\mathfrak{su}(2)$ representation is unitarily equivalent.

Define the dicyclic reflection character
\[
\chi(a^kx^e)=(-1)^e,
\]
and the corresponding diagonal involution
\[
\Gamma_\chi|a^kx^e\rangle=(-1)^e|a^kx^e\rangle.
\]
Since $\chi(a^6)=1$, $\Gamma_\chi$ commutes with the central $Z$ fiber.  For an orientation label $\sigma=\pm1$, set
\[
P_L=\frac{I-\sigma\Gamma_\chi}{2},
\qquad P_R=I-P_L.
\]
Finally define
\begin{equation}
\label{eq:weak}
T_i=\frac12P_L\tau_i.
\end{equation}

\begin{theorem}[Chiral weak representation]
The operators in \eqref{eq:weak} satisfy
\[
[T_i,T_j]=i\epsilon_{ijk}T_k,
\]
and vanish on $P_R\mathcal H_{504}$.  Thus $P_L\mathcal H_{504}$ carries weak doublets while $P_R\mathcal H_{504}$ carries weak singlets.
\end{theorem}
\begin{proof}
$P_L$ commutes with each $\tau_i$ and is a projector. Hence
\[
[T_i,T_j]=\frac14P_L[\tau_i,\tau_j]
=i\epsilon_{ijk}\frac12P_L\tau_k=i\epsilon_{ijk}T_k.
\]
Also $P_RT_i=0$.
\end{proof}

\section{Hypercharge and particle representation}
Set $T_3=P_LB/2$ and define $Y$ by the Gell--Mann--Nishijima relation
\[
Q=T_3+\frac Y2.
\]
Using \eqref{eq:Q},
\begin{equation}
\label{eq:Y}
Y=2(Q-T_3)=\frac{2A-I}{3}+P_RB.
\end{equation}

\begin{theorem}[Gauge quantum numbers]
After suppressing address multiplicities, the simultaneous $(A,B,P_L)$ sectors carry the following representations of
$SU(3)_c\times SU(2)_L\times U(1)_Y$:
\end{theorem}

\begin{center}
\begin{tabular}{@{}llllll@{}}
\toprule
sector & $A$ & $B$ & $SU(3)_c$ & $SU(2)_L$ & $Y$ \\
\midrule
$L_L=(\nu_L,e_L)$ & $-1$ & $\pm1$ & $\mathbf1$ & $\mathbf2$ & $-1$ \\
$e_R$ & $-1$ & $-1$ & $\mathbf1$ & $\mathbf1$ & $-2$ \\
$\nu_R$ & $-1$ & $+1$ & $\mathbf1$ & $\mathbf1$ & $0$ \\
$Q_L=(u_L,d_L)$ & $+1$ & $\pm1$ & $\mathbf3$ & $\mathbf2$ & $1/3$ \\
$d_R$ & $+1$ & $-1$ & $\mathbf3$ & $\mathbf1$ & $-2/3$ \\
$u_R$ & $+1$ & $+1$ & $\mathbf3$ & $\mathbf1$ & $4/3$ \\
\bottomrule
\end{tabular}
\end{center}

The electric charges are respectively
\[
Q(\nu)=0,\quad Q(e)=-1,\quad Q(u)=\frac23,\quad Q(d)=-\frac13.
\]
The independent factor $\mathbb C[\mathbb Z_3]$ supplies three family copies of this representation.

\begin{proof}
On $P_L$, equation \eqref{eq:Y} becomes
\[
Y_L=\frac{2A-I}{3},
\]
which is $-1$ for the center and $1/3$ for the leaves, independent of $B$ as required for a weak doublet.  On $P_R$,
\[
Y_R=\frac{2A-I}{3}+B.
\]
Substitution of $A=\pm1$ and $B=\pm1$ gives the four singlet hypercharges shown in the table.  The color statement follows from the previous theorem.
\end{proof}

\section{Address multiplicities}
The finite addressed representation contains additional factors on which the gauge representation acts trivially.  In particular, after fixing family, chirality and weak component, the lepton sectors have a six-dimensional residual address factor, while the quark sectors have
\[
\mathbb C^3_{\rm col}\otimes\mathbb C^2_{\rm ori}\otimes\mathbb C^6_{\rm addr}.
\]
The gauge representation type is obtained by suppressing these spectator address factors.  This distinction separates finite provenance multiplicity from particle-species multiplicity.

\section{Anomaly cancellation}
Use the left-handed Weyl convention.  One family, including a right-handed neutrino, is
\[
Q_L:(\mathbf3,\mathbf2)_{1/3},\qquad
L_L:(\mathbf1,\mathbf2)_{-1},
\]
\[
\nu_R^c:(\mathbf1,\mathbf1)_0,\qquad
e_R^c:(\mathbf1,\mathbf1)_2,
\]
\[
 u_R^c:(\overline{\mathbf3},\mathbf1)_{-4/3},\qquad
 d_R^c:(\overline{\mathbf3},\mathbf1)_{2/3}.
\]
The anomaly coefficients are
\[
\mathcal A_{SU(3)^2Y}=0,
\quad
\mathcal A_{SU(2)^2Y}=0,
\quad
\mathcal A_{\mathrm{grav}^2Y}=0,
\quad
\mathcal A_{Y^3}=0.
\]
There are four left-handed $SU(2)$ doublets per family after color multiplicity, so the global Witten anomaly is absent.

\section{Trace metric and one-loop running}
For one reduced family and generators normalized by $\operatorname{tr}(t_at_b)=\delta_{ab}/2$,
\[
G_3=\operatorname{Tr}(T_c^2)=2,
\qquad
G_2=\operatorname{Tr}(T_w^2)=2,
\qquad
G_Y=\operatorname{Tr}\left(\frac Y2\right)^2=\frac{10}{3}.
\]
Hence
\[
G_3:G_2:G_Y=1:1:\frac53.
\]
For a common parent kinetic normalization this gives
\[
g_3^2=g_2^2,
\qquad
\frac{g_Y^2}{g_2^2}=\frac35,
\qquad
\sin^2\theta_W=\frac38
\]
at the normalization scale.

With three families and one complex Higgs doublet of $Y_H=1$, the one-loop coefficients in
\[
\frac{dg_i}{d\ln\mu}=\frac{b_i}{16\pi^2}g_i^3
\]
are
\[
\boxed{b_3=-7,\qquad b_2=-\frac{19}{6},\qquad b_Y=\frac{41}{6}.}
\]
For $g_1=\sqrt{5/3}\,g_Y$, $b_1=41/10$.

\section{Computational verification}
The finite matrices were constructed directly in dimension $504$.  The following residuals were obtained in double precision:
\[
\|Q_{\rm polynomial}-Q_{\rm linear}\|_F=0,
\quad
\|Y-Y_{\rm closed}\|_F=0,
\]
\[
\max_{i,j}\|[T_i,T_j]-i\epsilon_{ijk}T_k\|_F=0,
\]
\[
\max_{a,i}\|[T_a^{(3)},T_i]\|_F=0,
\quad
\max_a\|[Y,T_a^{(3)}]\|_F=0,
\quad
\max_i\|[Y,T_i]\|_F=0.
\]
The color Casimir residual is $3.14\times10^{-16}$.  The complete machine-readable derivation and independent verifier accompany this manuscript.

\section{Conclusion}
The finite addressed Hilbert space $V_7\otimes\mathbb C[\mathbb Z_3\times\mathrm{Dic}_6]$ supports an explicit commuting representation of $\mathfrak{su}(3)_c\oplus\mathfrak{su}(2)_L\oplus\mathfrak u(1)_Y$.  The center/leaf involution determines lepton versus quark hypercharge, the central dicyclic involution supplies weak $T_3$, the reflection character supplies chirality, the six-leaf CRT factor supplies color, and the independent $\mathbb Z_3$ factor supplies three families.  The resulting reduced representation has the Standard-Model gauge quantum numbers with a right-handed neutrino and is anomaly free.

\begin{thebibliography}{9}
\bibitem{KogutSusskind1975} J. Kogut and L. Susskind, Hamiltonian formulation of Wilson's lattice gauge theories, \emph{Phys. Rev. D} \textbf{11}, 395 (1975).
\bibitem{Wiese2021} U.-J. Wiese, From quantum link models to D-theory: a resource efficient framework for the quantum simulation and computation of gauge theories, arXiv:2107.09335.
\bibitem{OsterwalderSeiler1978} K. Osterwalder and E. Seiler, Gauge field theories on a lattice, \emph{Ann. Phys.} \textbf{110}, 440--471 (1978).
\end{thebibliography}
\end{document}
